\frfilename{mt265.tex}
\versiondate{3.9.13}
\copyrightdate{2000}

\def\chaptername{Perubahan variabel dalam integral}
\def\sectionname{Ukuran permukaan}

\newsection{265}

Dalam bagian ini saya menyajikan versi baru argumen-argumen \S263, kali
ini bukan untuk membenarkan integrasi dengan substitusi, melainkan untuk
memberikan metode yang efektif dalam praktik guna menghitung ukuran
Hausdorff berdimensi $r$ dari permukaan mulus berdimensi $r$ dalam ruang
berdimensi $s$. Kasus dasar yang patut diingat adalah $r=2$, $s=3$,
meskipun kombinasi lain yang mudah Anda bayangkan juga akan sangat
membantu intuisi. Saya memberikan teorema fundamental (265E) yang
menyediakan rumus yang dapat digunakan untuk menghitung ukuran
berdimensi $r$ dari permukaan dalam ruang berdimensi $s$ yang
diparameterkan oleh fungsi terdiferensialkan, lalu mengerjakan beberapa
perhitungan untuk permukaan bola berdimensi $r$ (265F-265H). %265F 265G 265H

\leader{265A}{Ukuran Hausdorff ternormalisasi}\dvro{ Dalam}{ Seperti
yang saya kemukakan pada akhir bagian sebelumnya, ukuran Hausdorff,
sebagaimana didefinisikan dalam 264A-264C, bukanlah ukuran yang paling
sesuai untuk pekerjaan kita di sini; karena itu dalam}
bagian ini saya akan menggunakan {\bf ukuran Hausdorff ternormalisasi},
yang berarti $\nu_r=2^{-r}\beta_r\mu_{Hr}$, dengan $\mu_{Hr}$ ukuran
Hausdorff berdimensi $r$\cmmnt{ (ditafsirkan dalam ruang mana pun yang
sedang ditinjau)} dan $\beta_r=\mu_rB(\tbf{0},1)$ ukuran Lebesgue dari
setiap bola berjari-jari $1$ dalam $\BbbR^r$. \cmmnt{Akan memudahkan
jika kita mengambil $\beta_0=1$.} \dvro{Normalisasi ini}{Seperti
ditunjukkan dalam 264H-264I, normalisasi ini} membuat $\nu_r$ pada
$\BbbR^r$ sama dengan ukuran Lebesgue $\mu_r$.
\cmmnt{Perhatikan bahwa tentu saja}
$\nu_r^*=2^{-r}\beta_r\mu_{Hr}^*$\cmmnt{ (264Fb)}.

\leader{265B}{Subruang linear \dvrocolon{}}\cmmnt{ Sama seperti dalam
\S263, langkah pertama adalah menangani operator linear.

\medskip

\noindent}{\bf Teorema} Misalkan $r$, $s$ adalah bilangan bulat dengan
$1\le r\le s$, dan $T$ adalah matriks real berukuran $s\times r$;
pandang $T$ sebagai operator linear dari $\BbbR^r$ ke $\BbbR^s$.
Tetapkan $J=\sqrt{\det T\trs T}$, dengan $T\trs $ transpos dari $T$.
Tuliskan $\nu_r$ untuk ukuran Hausdorff ternormalisasi berdimensi $r$
pada $\BbbR^s$, $\Tau_r$ untuk domainnya, dan $\mu_r$ untuk ukuran
Lebesgue pada $\BbbR^r$. Maka

\Centerline{$\nu_r T[E]=J\mu_rE$}

\noindent untuk setiap himpunan terukur $E\subseteq\BbbR^r$. Jika $T$
injektif (yakni jika $J\ne 0$), maka

\Centerline{$\nu_rF=J\mu_rT^{-1}[F]$}

\noindent kapan pun $F\in\Tau_r$ dan $F\subseteq T[\BbbR^r]$.

\proof{ Rumus untuk $J$ mengandaikan bahwa $\det T\trs T$ nonnegatif,
yang belum kita ketahui; tetapi argumen berikut akan segera
membuktikannya dengan memadai.
\medskip
{\bf (a)} Misalkan $V$ adalah subruang linear dari $\BbbR^s$ yang
terdiri atas vektor $y=(\eta_1,\ldots,\eta_s)$ sedemikian sehingga
$\eta_i=0$ kapan pun $r<i\le s$. Misalkan $R$ adalah matriks
$r\times s$, $\langle\rho_{ij}\rangle_{i\le r,j\le s}$, dengan
$\rho_{ij}=1$ jika $i=j\le r$, dan $0$ jika tidak; maka matriks
$s\times r$ $R\trs $ dapat dipandang sebagai bijeksi dari $\BbbR^r$
ke $V$. Misalkan $W$ adalah subruang linear berdimensi $r$ dari
$\BbbR^s$ yang memuat $T[\BbbR^r]$, dan misalkan $P$ adalah matriks
ortogonal $s\times s$ sedemikian sehingga $P[W]=V$. Maka $S=RPT$
adalah matriks $r\times r$. Kita mempunyai $R\trs Ry=y$ untuk $y\in V$,
sehingga $R\trs RPT=PT$ dan

\Centerline{$S\trs S=T\trs P\trs R\trs RPT=T\trs P\trs PT=T\trs T$;}

\noindent jadi

\Centerline{$\det T\trs T=\det S\trs S=(\det S)^2\ge 0$}

\noindent dan $J=|\det S|$. Pada saat yang sama,

\Centerline{$P\trs R\trs S=P\trs R\trs RPT=P\trs PT=T$.}

\noindent Perhatikan bahwa $J=0$ jika dan hanya jika $S$ tidak
injektif, yakni jika dan hanya jika $T$ tidak injektif.

\medskip

{\bf (b)} Jika kita memandang matriks $s\times r$ $P\trs R\trs $ sebagai
pemetaan dari $\BbbR^r$ ke $\BbbR^s$, tampak bahwa
$\phi=P\trs R\trs $ adalah isometri antara $\BbbR^r$ dan $W$, dengan
invers $\phi^{-1}=RP\restrp W$. Karena itu $\phi$ adalah isomorfisme
antara ruang ukur $(\Bbb R^r,\mu_{Hr}^{(r)})$ dan
$(W,\mu_{HrW}^{(s)})$, dengan $\mu_{Hr}^{(r)}$ ukuran Hausdorff
berdimensi $r$ pada $\BbbR^r$ dan $\mu_{HrW}^{(s)}$ ukuran subruang
pada $W$ yang diinduksi oleh ukuran Hausdorff berdimensi $r$
$\mu_{Hr}^{(s)}$ pada $\BbbR^s$.

\medskip

\Prf\ {\bf (i)} Jika $A\subseteq\BbbR^r$ dan $A'\subseteq W$,

\Centerline{$\mu_{Hr}^{(s)*}(\phi[A])\le\mu_{Hr}^{(r)*}(A)$,
\quad $\mu_{Hr}^{(r)*}(\phi^{-1}[A'])\le\mu_{Hr}^{(s)*}(A')$,}

\noindent dengan menggunakan 264G dua kali. Jadi
$\mu_{Hr}^{(s)*}(\phi[A])=\mu_{Hr}^{(r)*}(A)$ untuk setiap
$A\subseteq \BbbR^r$.

\medskip

\quad{\bf (ii)} Sekarang, karena $W$ tertutup dan dengan demikian
berada dalam domain $\mu_{Hr}^{(s)}$ (264E), ukuran subruang
$\mu_{HrW}^{(s)}$ tepat merupakan ukuran yang diinduksi oleh
$\mu_{Hr}^{(s)*}\restrp W$ melalui metode \Caratheodory\
(214H(b-ii)). Karena $\phi$ adalah isomorfisme antara
$(\Bbb R^r,\mu_{Hr}^{(r)*})$ dan
$(W,\mu_{Hr}^{(s)*}\restrp W)$, ia merupakan isomorfisme antara
$(\BbbR^r,\mu_{Hr}^{(r)})$ dan $(W,\mu_{HrW}^{(s)})$.
\Qed

\medskip

{\bf (c)} Akibatnya, $\phi$ juga merupakan isomorfisme antara versi
ternormalisasi $(\BbbR^r,\mu_r)$ dan $(W,\nu_{rW})$, dengan
$\nu_{rW}$ ukuran subruang pada $W$ yang diinduksi oleh $\nu_r$.

Sekarang, jika $E\subseteq\BbbR^r$ terukur Lebesgue, berdasarkan
263A kita mempunyai $\mu_rS[E]=J\mu_rE$; sehingga

\Centerline{$\nu_rT[E]=\nu_r(P\trs R\trs [S[E]])=\nu_r(\phi[S[E]])
=\mu_rS[E]=J\mu_rE$.}

\noindent Jika $T$ injektif, maka $S=\phi^{-1}T$ juga harus injektif,
sehingga $J\ne 0$ dan

\Centerline{$\nu_rF=\mu_r(\phi^{-1}[F])
=J\mu_r(S^{-1}[\phi^{-1}[F]])=J\mu_rT^{-1}[F]$}

\noindent kapan pun $F\in\Tau_r$ dan $F\subseteq W=T[\BbbR^r]$.
}%end of proof of 265B

\leader{265C}{Akibat} Di bawah syarat-syarat 265B,

\Centerline{$\nu_r^*T[A]=J\mu_r^*A$}

\noindent untuk setiap $A\subseteq\BbbR^r$.

\proof{{\bf (a)} Jika $E$ terukur Lebesgue dan $A\subseteq E$,
maka $T[A]\subseteq T[E]$, sehingga

\Centerline{$\nu_r^*T[A]\le \nu_rT[E]=J\mu_rE$;}

\noindent karena $E$ sebarang, $\nu_r^*T[A]\le J\mu_r^*A$.

\medskip

{\bf (b)} Jika $J=0$, kita selesai. Jika $J\ne 0$, maka $T$ injektif,
sehingga jika $F\in\Tau_r$ dan $T[A]\subseteq F$, kita akan mempunyai

\Centerline{$J\mu_r^*A\le J\mu_rT^{-1}[F\cap T[\BbbR^r]]
=\nu_r(F\cap T[\BbbR^r])\le\nu_rF$;}

\noindent karena $F$ sebarang, $J\mu_r^*A\le\nu_r^*T[A]$.
}%end of proof of 265C

\leader{265D}{}\cmmnt{ Sekarang saya beralih ke lema yang bersesuaian
dengan 263C.

\medskip

\noindent}{\bf Lema} Misalkan $1\le r\le s$ dan $T$ adalah matriks
$s\times r$; tetapkan $J=\sqrt{\det T\trs T}$, dan misalkan $J\ne 0$.
Maka untuk setiap $\epsilon>0$ terdapat
$\zeta=\zeta(T,\epsilon)>0$ sedemikian sehingga

(i) $|\sqrt{\det S\trs S}-J|\le\epsilon$ kapan pun $S$ adalah matriks
$s\times r$ dan $\|S-T\|\le\zeta$\cmmnt{, dengan norma matriks
didefinisikan seperti dalam 262H};

(ii) kapan pun $D\subseteq\BbbR^r$ adalah himpunan terbatas dan
$\phi:D\to\Bbb R^s$ adalah fungsi sedemikian sehingga
$\|\phi(x)-\phi(y)-T(x-y)\|\le\zeta\|x-y\|$ untuk semua $x$, $y\in D$,
maka $|\nu_r^*\phi[D]-J\mu_r^*D|\le\epsilon\mu_r^*D$.

\proof{{\bf (a)} Karena $\det S\trs S$ adalah fungsi kontinu dari
koefisien-koefisien $S$, 262Hb memberi tahu kita bahwa harus ada
$\zeta_0>0$ sedemikian sehingga
$|J-\sqrt{\det S\trs S}|\le\epsilon$ kapan pun
$\|S-T\|\le\zeta_0$.

\medskip

{\bf (b)} Karena $J\ne 0$, $T$ injektif, dan terdapat matriks
$r\times s$ $T^*$ sedemikian sehingga $T^*T$ adalah matriks identitas
$r\times r$. Ambil $\zeta>0$ sedemikian sehingga
$\zeta\le\zeta_0$, $\zeta\|T^*\|<1$,
$J(1+\zeta\|T^*\|)^r\le J+\epsilon$, dan
$1-J^{-1}\epsilon\le(1-\zeta\|T^*\|)^r$.

Misalkan $\phi:D\to\BbbR^s$ sedemikian sehingga
$\|\phi(x)-\phi(y)-T(x-y)\|\le\zeta\|x-y\|$ kapan pun $x$, $y\in D$.
Tetapkan $\psi=\phi T^*$, sehingga $\phi=\psi T$. Maka

\Centerline{$\|\psi(u)-\psi(v)\|\le(1+\zeta\|T^*\|)\|u-v\|$,
\quad$\|u-v\|\le(1-\zeta\|T^*\|)^{-1}\|\psi(u)-\psi(v)\|$}

\noindent kapan pun $u$, $v\in T[D]$.
\Prf\ Ambil $x$, $y\in D$ sedemikian sehingga $u=Tx$, $v=Ty$;
tentu saja $x=T^*u$, $y=T^*v$. Maka

$$\eqalign{\|\psi(u)-\psi(v)\|
&=\|\phi(T^*u)-\phi(T^*v)\|
=\|\phi(x)-\phi(y)\|\cr
&\le\|T(x-y)\|+\zeta\|x-y\|\cr
&=\|u-v\|+\zeta\|T^*u-T^*v\|
\le\|u-v\|(1+\zeta\|T^*\|).\cr}$$

\noindent Selanjutnya,

$$\eqalign{\|u-v\|
&=\|Tx-Ty\|
\le\|\phi(x)-\phi(y)\|+\zeta\|x-y\|\cr
&=\|\psi(u)-\psi(v)\|+\zeta\|T^*u-T^*v\|\cr
&\le\|\psi(u)-\psi(v)\|+\zeta\|T^*\|\|u-v\|,\cr}$$

\noindent sehingga
$(1-\zeta\|T^*\|)\|u-v\|\le\|\psi(u)-\psi(v)\|$ dan
$\|u-v\|\le(1-\zeta\|T^*\|)^{-1}\|\psi(u)-\psi(v)\|$. \Qed

\medskip

{\bf (c)} Sekarang, dari 264G dan 265C kita melihat bahwa

\Centerline{$\nu^*_r\phi[D]
=\nu_r^*\psi[T[D]]
\le(1+\zeta\|T^*\|)^r\nu_r^*T[D]
=(1+\zeta\|T^*\|)^rJ\mu_r^*D
\le(J+\epsilon)\mu_r^*D$,}

\noindent dan (asalkan $\epsilon\le J$)

$$\eqalignno{(J-\epsilon)\mu_r^*D
&=(1-J^{-1}\epsilon)\nu_r^*T[D]
\le(1-J^{-1}\epsilon)(1-\zeta\|T^*\|)^{-r}\nu_r^*\psi[T[D]]\cr
\displaycause{dengan menerapkan 264G pada $\psi^{-1}:\psi[T[D]]\to T[D]$}
&\le\nu_r^*\psi[T[D]]
=\nu_r^*\phi[D].\cr}$$

\noindent (Tentu saja, jika $\epsilon\ge J$, maka pasti
$(J-\epsilon)\mu_r^*D\le\nu_r^*\phi[D]$.) Jadi
\Centerline{$(J-\epsilon)\mu_r^*D\le\nu_r^*\phi[D]
\le(J+\epsilon)\mu_r^*D$}

\noindent seperti yang diperlukan, dan kita mempunyai $\zeta$ yang
sesuai.
}%end of proof of 265D

\leader{265E}{Teorema} Misalkan $1\le r\le s$; tuliskan $\mu_r$ untuk
ukuran Lebesgue pada $\BbbR^r$, $\nu_r$ untuk ukuran Hausdorff
ternormalisasi pada $\BbbR^s$, dan $\Tau_r$ untuk domain $\nu_r$.
Misalkan $D\subseteq\BbbR^r$ sebarang, dan
$\phi:D\to\BbbR^s$ adalah fungsi yang terdiferensialkan relatif terhadap
domainnya pada setiap titik dalam $D$. Untuk setiap $x\in D$, misalkan $T(x)$
adalah suatu turunan dari $\phi$ di $x$ relatif terhadap $D$, dan
tetapkan $J(x)=\sqrt{\det T(x)\trs T(x)}$.
Tetapkan $D'=\{x:x\in D,\,J(x)>0\}$. Maka

\quad (i) $J:D\to\coint{0,\infty}$ adalah fungsi terukur;

\quad (ii) $\nu_r^*\phi[D]\le\int_DJ(x)\mu_r(dx)$,

\noindent dengan memperbolehkan $\infty$ sebagai nilai integral;

\quad (iii) $\nu_r^*\phi[D\setminus D']=0$.

\noindent Jika $D$ terukur Lebesgue, maka

\quad (iv) $\phi[D]\in\Tau_r$.

\noindent Jika $D$ terukur dan $\phi$ injektif, maka

\quad (v) $\nu_r\phi[D]=\int_DJ\,d\mu_r$;

\quad (vi) untuk setiap himpunan $E\subseteq\phi[D]$, $E\in\Tau_r$
jika dan hanya jika $\phi^{-1}[E]\cap D'$ terukur Lebesgue, dan
dalam hal ini

\Centerline{$\nu_rE
=\int_{\phi^{-1}[E]}J(x)\mu_r(dx)
=\int_DJ\times\chi(\phi^{-1}[E])d\mu_r$;}

\quad (vii) untuk setiap fungsi bernilai real $g$ yang didefinisikan
pada suatu himpunan bagian dari $\phi[D]$,

\Centerline{$\int_{\phi[D]}g\,d\nu_r=\int_DJ\times g\phi\,d\mu_r$}

\noindent jika salah satu integral terdefinisi dalam
$[-\infty,\infty]$, dengan ketentuan bahwa kita menafsirkan
$J(x)g(\phi(x))$ sebagai nol ketika $J(x)=0$ dan $g(\phi(x))$ tidak
terdefinisi.

\proof{ Saya akan mengikuti alur yang digariskan dalam bukti 263D.

\medskip

{\bf (a)} Sama seperti dalam 263D, kita tahu bahwa $J:D\to\Bbb R$
terukur, sebab $J(x)$ adalah fungsi kontinu dari koefisien-koefisien
$T(x)$, yang semuanya terukur berdasarkan 262P. Jika $D$ terukur
Lebesgue, terdapat barisan $\sequencen{F_n}$ himpunan bagian kompak dari
$D$ sedemikian sehingga $D\setminus\bigcup_{n\in\Bbb N}F_n$
terabaikan terhadap $\mu_r$. Sekarang $\phi[F_n]$ kompak, sehingga
termasuk dalam $\Tau_r$, untuk setiap $n\in\Bbb N$. Adapun
$\phi[D\setminus\bigcup_{n\in\Bbb N}F_n]$, himpunan ini harus
terabaikan terhadap $\nu_r$ berdasarkan 264G, karena $\phi$ merupakan
gabungan terhitung dari fungsi-fungsi Lipschitz (262N). Jadi

\Centerline{$\phi[D]=\bigcup_{n\in\Bbb N}
  \phi[F_n]\cup\phi[D\setminus\bigcup_{n\in\Bbb N}F_n]\in\Tau_r$.}

\noindent Ini menangani (i) dan (iv).

\medskip

{\bf (b)} Untuk sementara, andaikan $D$ terbatas dan $J(x)>0$ untuk
setiap $x\in D$, serta tetapkan $\epsilon>0$. Misalkan $M^*_{sr}$
adalah himpunan matriks $s\times r$ $T$ sedemikian sehingga
$\det T\trs T\ne 0$, yakni pemetaan yang bersesuaian
$T:\BbbR^r\to\BbbR^s$ injektif. Untuk $T\in M_{sr}^*$, ambil
$\zeta(T,\epsilon)>0$ seperti dalam 265D.

Ambil $\sequencen{D_n}$, $\sequencen{T_n}$ seperti dalam 262M, dengan
$A=M^*_{sr}$, sehingga $\sequencen{D_n}$ adalah partisi $D$ menjadi
himpunan-himpunan yang relatif terukur dalam $D$, dan setiap $T_n$
adalah matriks $s\times r$ sedemikian sehingga

\Centerline{$\|T(x)-T_n\|\le\zeta(T_n,\epsilon)$ kapan pun $x\in D_n$,}

\Centerline{$\|\phi(x)-\phi(y)-T_n(x-y)\|\le\zeta(T_n,\epsilon)\|x-y\|$
untuk semua $x$, $y\in D_n$.}

\noindent Kemudian, dengan menetapkan
$J_n=\sqrt{\det T_n\trs T_n}$, kita mempunyai

\Centerline{$|J(x)-J_n|\le\epsilon$ untuk setiap $x\in D_n$,}

\Centerline{$|\nu_r^*\phi[D_n]-J_n\mu_r^*D_n|\le\epsilon\mu_r^*D_n$,}

\noindent berdasarkan pemilihan $\zeta(T_n,\epsilon)$. Jadi

$$\eqalignno{\nu_r^*\phi[D]
&\le\sum_{n=0}^{\infty}\nu_r^*\phi[D_n]\cr
\displaycause{karena $\phi[D]=\bigcup_{n\in\Bbb N}\phi[D_n]$}
&\le\sum_{n=0}^{\infty}(J_n\mu_r^*D_n+\epsilon\mu_r^*D_n)
\le\epsilon\mu_r^*D+\sum_{n=0}^{\infty}J_n\mu_r^*D_n\cr
\displaycause{karena $D_n$ saling lepas dan relatif terukur dalam $D$}
&=\epsilon\mu_r^*D+\int_D\sum_{n=0}^{\infty}J_n\chi D_nd\mu_r\cr
&\le\epsilon\mu_r^*D+\int_DJ(x)+\epsilon\mu_r(dx)
=2\epsilon\mu_r^*D+\int_DJ\,d\mu_r.\cr}$$

Jika $D$ terukur dan $\phi$ injektif, maka semua $D_n$ adalah himpunan
bagian $\BbbR^r$ yang terukur Lebesgue, sehingga semua
$\phi[D_n]$ diukur oleh $\nu_r$, dan himpunan-himpunan itu juga saling
lepas. Oleh karena itu

$$\eqalign{\int_{D}J\,d\mu_r
&\le\sum_{n=0}^{\infty}J_n\mu_r D_n+\epsilon\mu_r D\cr
&\le\sum_{n=0}^{\infty}(\nu_r\phi[D_n]
   +\epsilon\mu_r D_n)+\epsilon\mu_rD
=\nu_r\phi[D]+2\epsilon\mu_rD.\cr}$$

\noindent Karena $\epsilon$ sebarang, kita memperoleh

\Centerline{$\nu_r^*\phi[D]\le\int_DJ\,d\mu_r$,}

\noindent dan jika $D$ terukur serta $\phi$ injektif,

\Centerline{$\int_DJ\,d\mu_r\le\nu_r\phi[D]$;}

\noindent jadi kita memperoleh (ii) dan (v), dengan asumsi bahwa $D$
terbatas dan $J>0$ pada setiap titik dalam $D$.

\medskip

{\bf (c)} Sama seperti dalam bukti 263D, sekarang kita dapat melepaskan
asumsi bahwa $D$ terbatas dengan meninjau
$B_k=B(\tbf{0},k)\subseteq\BbbR^r$; asalkan $J>0$ pada setiap titik dalam $D$,
kita memperoleh

\Centerline{$\nu_r^*\phi[D]
=\lim_{k\to\infty}\nu_r^*\phi[D\cap B_k]
\le\lim_{k\to\infty}\int_{D\cap B_k}J\,d\mu_r
=\int_DJ\,d\mu_r$,}

\noindent dengan kesamaan jika $D$ terukur dan $\phi$ injektif.

\medskip

{\bf (d)} Sekarang kita mendapati bahwa
$\nu_r^*\phi[D\setminus D']=0$.

\medskip

\Prf\grheada\ Misalkan $\eta\in\ocint{0,1}$. Definisikan
$\psi_{\eta}:D\to\BbbR^{s+r}$ dengan menetapkan
$\psi_{\eta}(x)=(\phi(x),\eta x)$, dan identifikasikan $\BbbR^{s+r}$
dengan $\BbbR^s\times\BbbR^r$. $\psi_{\eta}$ terdiferensialkan relatif
terhadap domainnya pada setiap titik dalam $D$, dengan turunan
$\tilde T_{\eta}(x)$, yaitu matriks $(s+r)\times r$ yang $s$ baris
atasnya terdiri atas matriks $s\times r$ $T(x)$, dan $r$ baris bawahnya
adalah $\eta I_r$, dengan $I_r$ matriks identitas $r\times r$.
(Gunakan 262Ib.) Tentu saja $\tilde T_{\eta}(x)$, dipandang sebagai
pemetaan dari $\BbbR^r$ ke $\BbbR^{s+r}$, injektif, sehingga

\Centerline{$\tilde J_{\eta}(x)
=\sqrt{\det\tilde T_{\eta}(x)\trs \tilde T_{\eta}(x)}
=\sqrt{\det(T(x)\trs T(x)+\eta^2I)}>0$.}

\noindent Kita mempunyai
$\lim_{\eta\downarrow 0}\tilde J_{\eta}(x)=J(x)=0$ untuk
$x\in D\setminus D'$.

\medskip

\quad\grheadb\ Nyatakan $T(x)$ sebagai
$\langle\tau_{ij}(x)\rangle_{i\le s,j\le r}$ untuk setiap $x\in D$.
Tetapkan

\Centerline{$C_m=\{x:x\in D,\,\|x\|\le m,\,|\tau_{ij}(x)|\le m$ untuk
semua $i\le s,\,j\le r\}$}

\noindent untuk setiap $m\ge 1$. Untuk $x\in C_m$, semua koefisien
$\tilde T_{\eta}(x)$ mempunyai modulus paling besar $m$; akibatnya
(dengan menggunakan ketaksamaan paling kasar yang segera tersedia),
semua koefisien $\tilde T_{\eta}(x)\trs \tilde T_{\eta}(x)$ mempunyai
modulus paling besar $(r+s)m^2$ dan
$\tilde J_{\eta}(x)\le\sqrt{r!(s+r)^r}m^r$. Karena itu kita dapat
menggunakan Teorema Konvergensi Terdominasi Lebesgue untuk melihat bahwa

\Centerline{$\lim_{\eta\downarrow 0}\int_{C_m\setminus D'}\tilde
J_{\eta}d\mu_r=0$.}

\medskip

\quad\grheadc\ Misalkan $\tilde\nu_r$ adalah ukuran Hausdorff
ternormalisasi berdimensi $r$ pada $\BbbR^{s+r}$. Dengan menerapkan
bagian (b) bukti ini pada $\psi_{\eta}\restr C_m\setminus D'$, kita
melihat bahwa

\Centerline{$\tilde\nu_r^*\psi_{\eta}[C_m\setminus D']
\le\int_{C_m\setminus D'}\tilde J_{\eta}d\mu_r$.}

\noindent Sekarang kita mempunyai pemetaan alami
$P:\BbbR^{s+r}\to\BbbR^s$ yang diberikan dengan menetapkan
$P(\xi_1,\ldots,\xi_{s+r})=(\xi_1,\ldots,\xi_s)$, dan $P$ bersifat
$1$-Lipschitz, sehingga dengan 264G sekali lagi kita mempunyai (dengan
memperhitungkan konstanta normalisasi $2^{-r}\beta_r$)

\Centerline{$\nu_r^*P[A]\le\tilde\nu_r^*A$}

\noindent untuk setiap $A\subseteq\BbbR^{s+r}$. Khususnya,

\Centerline{$\nu_r^*\phi[C_m\setminus D']
=\nu_r^*P[\psi_{\eta}[C_m\setminus D']]
\le\tilde\nu_r^*\psi_{\eta}[C_m\setminus D']
\le\int_{C_m\setminus D'}\tilde J_{\eta}d\mu_r
\to 0$}

\noindent ketika $\eta\downarrow 0$. Namun ini berarti
$\nu_r^*\phi[C_m\setminus D']=0$. Karena
$D=\bigcup_{m\ge 1}C_m$, $\nu_r^*\phi[D\setminus D']=0$, seperti yang
diklaim. \Qed

\medskip

{\bf (e)} Ini membuktikan (iii) dalam teorema. Tentu saja ini cukup
untuk memberikan (ii) dan (v), sebab (dengan menerapkan (b)-(c) pada
$\phi\restr D'$) kita harus mempunyai

\Centerline{$\nu_r^*\phi[D]
=\nu_r^*\phi[D']
\le\int_{D'}J\,d\mu_r
=\int_DJ\,d\mu_r$,}

\noindent dengan kesamaan jika $D$ (dan karena itu juga $D'$) terukur
dan $\phi$ injektif.

\medskip

{\bf (f)} Sekarang mari kita beralih ke bagian (vi). Andaikan $D$
terukur dan $\phi$ injektif.

\medskip

\quad\grheada\ Misalkan $E\subseteq\phi[D]$ termasuk dalam $\Tau_r$.
Tetapkan

\Centerline{$H_k=\{x:x\in D,\,\|x\|\le k,\,J(x)\le k\}$}

\noindent untuk setiap $k$; maka setiap $H_k$ terukur Lebesgue,
sehingga (dengan menerapkan (iv) pada $\phi\restr H_k$)
$\phi[H_k]\in\Tau_r$, dan

\Centerline{$\nu_r\phi[H_k]\le k\mu_rH_k<\infty$.}

\noindent Jadi $\phi[D]$ dapat ditutupi oleh suatu barisan himpunan
berukuran berhingga terhadap $\nu_r$, yang tentu saja juga berukuran
berhingga terhadap ukuran Hausdorff berdimensi $r$ pada $\BbbR^s$.
Berdasarkan 264Fc, terdapat himpunan Borel
$E_1$, $E_2\subseteq\BbbR^s$ sedemikian sehingga
$E_1\subseteq E\subseteq E_2$ dan $\nu_r(E_2\setminus E_1)=0$.
Sekarang $F_1=\phi^{-1}[E_1]$, $F_2=\phi^{-1}[E_2]$ adalah himpunan
bagian terukur Lebesgue dari $D$, dan

\Centerline{$\int_{F_2\setminus F_1}J\,d\mu_r=\nu_r\phi[F_2\setminus
F_1]=\nu_r(\phi[D]\cap E_2\setminus E_1)=0$.}

\noindent Karena itu $\mu_r(D'\cap(F_2\setminus F_1))=0$. Namun karena

\Centerline{$D'\cap F_1\subseteq D'\cap\phi^{-1}[E]
\subseteq D'\cap F_2$,}

\noindent diperoleh bahwa $D'\cap\phi^{-1}[E]$ terukur, dan bahwa

$$\eqalign{\int_{\phi^{-1}[E]}J\,d\mu_r
&=\int_{D'\cap\phi^{-1}[E]}J\,d\mu_r
=\int_{D'\cap F_1}J\,d\mu_r\cr
&=\int_{D\cap F_1}J\,d\mu_r
=\nu_r\phi[D\cap F_1]
=\nu_rE_1
=\nu_rE.\cr}$$

\noindent Selain itu, $J\times\chi(\phi^{-1}[E])
=J\times\chi(D'\cap\phi^{-1}[E])$ terukur, sehingga kita dapat menulis
$\int J\times\chi(\phi^{-1}[E])$ sebagai pengganti
$\int_{\phi^{-1}[E]}J$.

\medskip

\quad\grheadb\ Jika $E\subseteq\phi[D]$ dan
$D'\cap\phi^{-1}[E]$ terukur, maka tentu saja

\Centerline{$E=\phi[D'\cap\phi^{-1}[E]]
  \cup\phi[(D\setminus D')\cap\phi^{-1}[E]]\in\Tau_r$,}

\noindent karena $\phi[G]\in\Tau_r$ untuk setiap himpunan terukur
$G\subseteq D$ dan $\phi[D\setminus D']$ terabaikan terhadap
$\nu_r$.

\medskip

{\bf (g)} Akhirnya, (vii) segera mengikuti dari (vi), dengan
menerapkan 235J pada $\mu_r$ dan ukuran subruang yang diinduksi oleh
$\nu_r$ pada $\phi[D]$.
}%end of proof of 265E

\leader{265F}{Permukaan bola}\cmmnt{ Untuk menunjukkan bagaimana
gagasan-gagasan ini dapat diterapkan pada salah satu kasus dasar, saya
memberikan rincian metode untuk mendeskripsikan ukuran permukaan bola
dalam ruang berdimensi $s$. Ambil $r\ge 1$ dan $s=r+1$.} Tuliskan
$S_r$ untuk $\{z:z\in\Bbb R^{r+1},\,\|z\|=1\}$, permukaan bola
berdimensi $r$.
\dvro{Maka ukuran Hausdorff ternormalisasi berdimensi $r$ dari $S_r$
adalah $2\pi\beta_{r-1}$, dengan $\beta_{r-1}$ volume bola satuan dari
$\BbbR^{r-1}$ (dengan menafsirkan $\beta_0$ sebagai $1$).}
{Maka kita mempunyai parametrisasi $\phi_r$ dari $S_r$ yang diberikan
dengan menetapkan

$$\phi_r\Matrix{\xi_1\\ \xi_2\\ \ldots \\ \ldots \\ \ldots\\ \xi_r}
=\Matrix{\sin\xi_1\sin\xi_2\sin\xi_3\ldots\sin\xi_r\\
\cos\xi_1\sin\xi_2\sin\xi_3\ldots\sin\xi_r\\
\cos\xi_2\sin\xi_3\ldots\sin\xi_r\\
\ldots\\
\cos\xi_{r-2}\sin\xi_{r-1}\sin\xi_r\\
\cos\xi_{r-1}\sin\xi_r\\
\cos\xi_r}.$$

\noindent Saya memilih perumusan ini karena saya ingin menggunakan
argumen induktif yang didasarkan pada kenyataan bahwa

\Centerline{$\phi_{r+1}\Matrix{x\\ \xi}
=\Matrix{\sin\xi\,\phi_r(x)\\ \cos\xi}$}
\noindent untuk $x\in\BbbR^r$, $\xi\in\Bbb R$. Setiap $\phi_r$
terdiferensialkan, berdasarkan 262Id. Jika kita menetapkan

$$\eqalign{D_r
=\{x:\xi_1&\in\ocint{-\pi,\pi},\,\xi_2,\ldots,\xi_r\in[0,\pi],\cr
&\text{ jika }\xi_j\in\{0,\pi\}\text{ maka }
  \xi_i=0\text{ untuk }i<j\},\cr}$$

\noindent maka mudah diperiksa bahwa $D_r$ adalah himpunan bagian Borel
dari $\Bbb R^r$ dan $\phi_r\restr D_r$ adalah bijeksi antara $D_r$ dan
$S_r$. Sekarang misalkan $T_r(x)$ adalah matriks $(r+1)\times r$
$\phi_r'(x)$. Maka

$$T_{r+1}\Matrix{x\\ \xi}=\Matrix{\sin\xi\,T_r(x)&\cos\xi\,\phi_r(x)\\
                                  \tbf{0}&-\sin\xi}.$$

\noindent Jadi

$$(T_{r+1}\Matrix{x\\ \xi})\trs T_{r+1}\Matrix{x\\ \xi}
=\Matrix{\sin^2\xi\,T_r(x)\trs T_r(x)&\sin\xi\cos\xi\,T_r(x)\trs \phi_r(x)\\
\cos\xi\sin\xi\,\phi_r(x)\trs T_r(x)&\cos^2\xi\phi_r(x)\trs \phi_r(x)+\sin^2\xi}
.$$

\noindent Namun tentu saja
$\phi_r(x)\trs \phi_r(x)=\|\phi_r(x)\|^2=1$ untuk setiap $x$, dan
(dengan mendiferensialkan terhadap setiap koordinat $x$, jika Anda
ingin) $T_r(x)\trs \phi_r(x)=\tbf{0}$,
$\phi_r(x)\trs T_r(x)=\tbf{0}$. Jadi kita memperoleh

$$(T_{r+1}\Matrix{x\\ \xi})\trs T_{r+1}\Matrix{x\\ \xi}
=\Matrix{\sin^2\xi\,T_r(x)\trs T_r(x)&\tbf{0}\\
\tbf{0}&1},$$

\noindent dan, dengan menulis
$J_r(x)=\sqrt{\det T_r(x)\trs T_r(x)}$,

$$J_{r+1}\Matrix{x\\ \xi}=|\sin^r\xi|J_r(x).$$

Pada titik ini kita melakukan induksi pada $r$ untuk melihat bahwa

\Centerline{$J_r(x)
=|\sin^{r-1}\xi_r\sin^{r-2}\xi_{r-1}\ldots\sin\xi_2|$}

\noindent (karena tentu saja induksi dimulai dengan kasus $r=1$,
\discrcenter{468pt}{$\phi_1(x)=\Matrix{\sin x\\ \cos x}$,
\ifdim\pagewidth>467pt\quad\fi
$T_1(x)=\Matrix{\cos x\\ -\sin x}$,
\ifdim\pagewidth>467pt\quad\fi
$T_1(x)\trs T_1(x)=1$,
\ifdim\pagewidth>467pt\quad\fi
$J_1(x)=1$).}

Untuk mencari ukuran permukaan $S_r$, kita perlu menghitung

$$\eqalignno{\int_{D_r}J_rd\mu_r
&=\int_0^{\pi}\ldots\int_0^{\pi}\int_{-\pi}^{\pi}
\sin^{r-1}\xi_r\ldots\sin\xi_2
d\xi_1d\xi_2\ldots d\xi_r\cr
&=2\pi\prod_{k=2}^r\int_0^{\pi}\sin^{k-1}t\,dt
=2\pi\prod_{k=1}^{r-1}\int_{-\pi/2}^{\pi/2}\cos^kt\,dt\cr}$$

\noindent (dengan menyubstitusikan $\bover{\pi}{2}-t$ untuk $t$).
Namun, dalam bahasa 252Q, ini tepat sama dengan

\Centerline{$2\pi\prod_{k=1}^{r-1}I_k=2\pi\beta_{r-1}$,}

\noindent dengan $\beta_{r-1}$ volume bola satuan dari
$\BbbR^{r-1}$\cmmnt{ (dengan menafsirkan $\beta_0$ sebagai $1$, jika
Anda mau)}.
}%end of dvro

\leader{265G}{}\cmmnt{ Luas permukaan bola juga dapat dihitung melalui
hasil berikut.

\medskip

\noindent}{\bf Teorema} Misalkan $\mu_{r+1}$ adalah ukuran Lebesgue
pada $\BbbR^{r+1}$, dan $\nu_r$ ukuran Hausdorff ternormalisasi
berdimensi $r$ pada $\BbbR^{r+1}$. Jika $f$ adalah fungsi bernilai real
yang terintegralkan secara lokal terhadap $\mu_{r+1}$,
$y\in\BbbR^{r+1}$, dan $\delta>0$,

\Centerline{$\int_{B(y,\delta)}fd\mu_{r+1}
=\int_0^{\delta}\int_{\partial B(y,t)}fd\nu_rdt$,}

\noindent dengan $\partial B(y,t)$ menyatakan permukaan bola
$\{x:\|x-y\|=t\}$\cmmnt{ dan integral $\int\ldots dt$ diambil terhadap
ukuran Lebesgue pada $\Bbb R$}.

\proof{ Ambil sembarang fungsi terdiferensialkan
$\phi:\BbbR^r\to S_r$ bersama himpunan Borel $F\subseteq\BbbR^r$
sedemikian sehingga $\phi\restr F$ merupakan bijeksi antara $F$ dan
$S_r$; pasangan $(\phi,F)$ semacam itu dideskripsikan dalam 265F.
Definisikan $\psi:\BbbR^r\times\Bbb R\to\BbbR^{r+1}$ dengan
menetapkan $\psi(z,t)=y+t\phi(z)$; maka $\psi$ terdiferensialkan dan
$\psi\restr F\times\ocint{0,\delta}$ merupakan bijeksi antara
$F\times\ocint{0,\delta}$ dan $B(y,\delta)\setminus\{y\}$. Untuk
$t\in\ocint{0,\delta}$, $z\in\BbbR^r$, tetapkan
$\psi_t(z)=\psi(z,t)$; maka $\psi_t\restr F$ merupakan bijeksi antara
$F$ dan permukaan bola $\{x:\|x-y\|=t\}=\partial B(y,t)$.

Turunan $\phi$ di $z$ adalah matriks $(r+1)\times r$, katakanlah
$T_1(z)$, dan turunan $T_t(z)$ dari $\psi_t$ di $z$ tepat $tT_1(z)$;
turunan $\psi$ di $(z,t)$ adalah matriks $(r+1)\times(r+1)$
$T(z,t)=\Matrix{tT_1(z)&\phi(z)}$, dengan $\phi(z)$ ditafsirkan sebagai
vektor kolom. Jika kita menetapkan

\Centerline{$J_t(z)=\sqrt{\det T_t(z)\trs T_t(z)}$,
\quad$J(z,t)=|\det T(z,t)|$,}

\noindent maka

$$\eqalign{J(z,t)^2
=\det T(z,t)\trs T(z,t)
&=\det\Matrix{tT_1(z)\trs \\ \phi(z)\trs }\Matrix{tT_1(z)&\phi(z)}\cr
&=\det\Matrix{t^2T_1(z)\trs T_1(z)&\tbf{0}\\ \tbf{0}&1}
=J_t(z)^2,\cr}$$

\noindent sebab ketika kita menghitung koefisien $(i,r+1)$ dari
$T(z,t)\trs T(z,t)$, untuk $1\le i\le r$, koefisien itu adalah

\Centerline{$\sum_{j=1}^{r+1}t\Pd{\phi_j}{\zeta_i}(z)\phi_j(z)
=\Bover{t}2\Pd{}{\zeta_i}(\sum_{j=1}^{r+1}\phi_j(z)^2)
=0$,}

\noindent dengan $\phi_j$ koordinat ke-$j$ dari $\phi$; sedangkan
koefisien $(r+1,r+1)$ dari $T(z,t)\trs T(z,t)$ tepat
$\sum_{j=1}^{r+1}\phi_j(z)^2=1$. Jadi sebenarnya
$J(z,t)=J_t(z)$ untuk semua $z\in\BbbR^r$, $t>0$.

Sekarang, untuk $f\in\eusm L^1(\mu_{r+1})$, kita dapat menghitung

$$\eqalignno{\int_{B(y,\delta)}fd\mu_{r+1}
&=\int_{B(y,\delta)\setminus\{y\}}fd\mu_{r+1}\cr
&=\int_{F\times\ocint{0,\delta}}f(\psi(z,t))J(z,t)\mu_{r+1}(d(z,t))\cr
\displaycause{berdasarkan 263D}
&=\int_0^{\delta}\int_Ff(\psi_t(z))J_t(z)\mu_r(dz)dt\cr
\displaycause{dengan $\mu_r$ ukuran Lebesgue pada $\BbbR^r$,
berdasarkan teorema Fubini, 252B}
&=\int_0^{\delta}\int_{\partial B(y,t)}fd\nu_rdt\cr}$$

\noindent berdasarkan 265E(vii).
}%end of proof of 265G

\leader{265H}{Akibat} Jika $\nu_r$ adalah ukuran Hausdorff
ternormalisasi berdimensi $r$ pada $\BbbR^{r+1}$, maka
$\nu_rS_r=(r+1)\beta_{r+1}$.

\proof{ Dalam 265G, ambil $y=\tbf{0}$, $\delta=1$, dan
$f=\chi B(\tbf{0},1)$; maka

\Centerline{$\beta_{r+1}=\int fd\mu_{r+1}
=\int_0^1\nu_r(\partial B(\tbf{0},t))dt=\int_0^1t^r\nu_rS_rdt
=\Bover1{r+1}\nu_rS_r$,}

\noindent kali ini dengan menerapkan 264G pada pemetaan
$x\mapsto tx$, $x\mapsto\bover1tx$ dari $\BbbR^{r+1}$ ke dirinya
sendiri untuk melihat bahwa
$\nu_r(\partial B(\tbf{0},t))=t^r\nu_rS_r$ untuk $t>0$.
}%end of proof of 265H

\exercises{
\leader{265X}{Latihan dasar (a)} Misalkan $r\ge 1$, dan misalkan
$S_r(\alpha)=\{z:z\in\BbbR^{r+1},\,\|z\|=\alpha\}$ adalah permukaan bola
berdimensi $r$ dan berjari-jari $\alpha$. Tunjukkan bahwa
$\nu_rS_r(\alpha)=2\pi\beta_{r-1}\alpha^r=(r+1)\beta_{r+1}\alpha^r$
untuk setiap $\alpha\ge 0$.

\sqheader 265Xb Misalkan $r\ge 1$, dan untuk $a\in[-1,1]$ tetapkan
$C_a=\{z:z\in\BbbR^{r+1},\,\|z\|=1,\,\zeta_{r+1}\ge a\}$, dengan
$z=(\zeta_1,\ldots,\zeta_{r+1})$ seperti biasa. (i) Tunjukkan bahwa

\Centerline{$\nu_rC_a=r\beta_r\int_0^{\arccos a}\sin^{r-1}t\,dt$.}

\noindent (ii)\dvAnew{2013}
Hitung integral tersebut untuk kasus $r=2$, $r=4$.

\sqheader 265Xc Tuliskan lagi
$C_a=\{z:z\in S_r,\,\zeta_{r+1}\ge a\}$, dengan
$S_r\subseteq\BbbR^{r+1}$ permukaan bola satuan. Tunjukkan bahwa, untuk setiap
$a\in\ocint{0,1}$,
$\nu_rC_a\le\Bover{\nu_rS_r}{2(r+1)a^2}$. \Hint{hitung
$\sum_{i=1}^{r+1}\int_{S_r}\|\xi_i\|^2\nu_r(dx)$.}
%265F

\sqheader 265Xd Misalkan $\phi:\ooint{0,1}\to\BbbR^r$ adalah fungsi
injektif yang terdiferensialkan. Tunjukkan bahwa `panjang' atau ukuran
Hausdorff satu-dimensi dari $\phi[\,\ooint{0,1}\,]$ tepat sama dengan
$\int_0^1\|\phi'(t)\|dt$.

\spheader 265Xe(i) Tunjukkan bahwa jika $I$ adalah matriks identitas
$r\times r$ dan $z\in\BbbR^r$, maka
$\det(I+zz\trs )=1+\|z\|^2$. \Hint{lakukan induksi pada $r$.}
(ii) Tuliskan $U_{r-1}$ untuk bola satuan terbuka dalam $\BbbR^{r-1}$,
dengan $r\ge 2$. Definisikan
$\phi:U_{r-1}\times\Bbb R\to S_r$ dengan menetapkan

$$\phi\Matrix{x\\ \xi}=\Matrix{x\\ \theta(x)\cos\xi \\
\theta(x)\sin\xi},$$

\noindent dengan $\theta(x)=\sqrt{1-\|x\|^2}$. Tunjukkan bahwa

$$\phi'\Matrix{x\\ \xi}\trs \phi'\Matrix{x\\ \xi}
=\Matrix{I+\Bover1{\theta(x)^2}xx\trs  & \tbf{0}\\
        \tbf{0} & \theta(x)^2},$$

\noindent sehingga $J\Matrix{x\\ \xi}=1$ untuk semua
$x\in U_{r-1}$, $\xi\in\Bbb R$. (iii) Karena itu, tunjukkan bahwa
ukuran Hausdorff ternormalisasi berdimensi $r$ dari
$\{y:y\in S_r,\,\sum_{i=1}^{r-1}\eta_i^2<1\}$ tepat
$2\pi\beta_{r-1}$, dengan $\beta_{r-1}$ ukuran Lebesgue dari
$U_{r-1}$. (iv) Dengan meninjau
$\psi z=\Matrix{z\\ 0\\ 0}$ untuk $z\in S_{r-2}$, atau dengan cara
lain, tunjukkan bahwa ukuran Hausdorff ternormalisasi berdimensi $r$
dari $S_r$ adalah $2\pi\beta_{r-1}$. (v) Kali ini dengan menetapkan
$C_a=\{z:z\in\BbbR^{r+1},\,\|z\|=1,\,\zeta_1\ge a\}$, tunjukkan bahwa
$\nu_rC_a=2\pi\mu_{r-1}\{x:x\in\BbbR^{r-1},\,\|x\|\le 1,\,\xi_1\ge a\}$
untuk setiap $a\in[-1,1]$.
%265F

\spheader 265Xf\dvAnew{2013} Misalkan $r\ge 2$. Dengan
mengidentifikasikan $\BbbR^r$ dengan $\BbbR^{r-1}\times\Bbb R$,
misalkan $C_r$ adalah silinder
$B_{r-1}\times[-1,1]\supseteq B_r$, dan
$\partial C_r=(B_{r-1}\times\{-1,1\})\cup(S_{r-2}\times[-1,1])$
batasnya. Tunjukkan bahwa

\Centerline{$\Bover{\mu_rB_r}{\mu_rC_r}
=\Bover{\nu_{r-1}S_{r-1}}{\nu_{r-1}(\partial C_r)}$.}

\noindent(Kasus $r=3$ berasal dari Archimedes.)
%265F

\leader{265Y}{Latihan lanjutan (a)}
%\spheader 265Ya
Ambil $a<b$ dalam $\Bbb R$. (i) Tunjukkan bahwa
$\phi:[a,b]\to\BbbR^r$ kontinu mutlak dalam pengertian 264Yp jika dan
hanya jika semua koordinatnya $\phi_i:[a,b]\to\Bbb R$, untuk $i\le r$,
kontinu mutlak dalam pengertian \S225. (ii) Misalkan
$\phi:[a,b]\to\BbbR^r$ adalah fungsi kontinu, dan tetapkan
$F=\{x:x\in\ooint{a,b},\,\phi$ terdiferensialkan di $x\}$. Tunjukkan
bahwa $\phi$ kontinu mutlak jika dan hanya jika
$\int_F\|\phi'(x)\|dx$ berhingga dan
$\nu_1(\phi[[a,b]\setminus F])=0$, dengan $\nu_1$ ukuran Hausdorff
satu-dimensi (ternormalisasi) pada $\BbbR^r$. ({\it Petunjuk\/}:
225K.) (iii) Tunjukkan bahwa jika $\phi:[a,b]\to\BbbR^r$ kontinu
mutlak, maka $\nu_1^*(\phi[D])\le\int_D\|\phi'(x)\|dx$ untuk setiap
$D\subseteq[a,b]$, dengan kesamaan jika $D$ terukur dan
$\phi\restr D$ injektif.

\spheader 265Yb Misalkan $a\le b$ dalam $\Bbb R$, dan
$f:[a,b]\to\Bbb R$ adalah fungsi kontinu bervariasi terbatas dengan
grafik $\Gamma_f$. Tunjukkan bahwa ukuran Hausdorff satu-dimensi dari
$\Gamma_f$ adalah
$\Var_{[a,b]}(f)+\int_a^b(\sqrt{1+(f')^2}-|f'|)$.
%mt26bits
}%end of exercises

\endnotes{
\Notesheader{265} Bukti 265B tampaknya menggunakan hampir seluruh
paruh kedua alfabet. Gagasannya seharusnya cukup langsung. Karena
$T[\BbbR^r]$ berdimensi paling besar $r$, himpunan itu dapat dirotasikan
oleh transformasi ortogonal $P$ ke dalam suatu subruang dari subruang
kanonik berdimensi $r$, $V$, yang merupakan salinan alami dari
$\BbbR^r$; matriks $R$ merepresentasikan proses penyalinan dari $V$ ke
$\BbbR^r$, dan $\phi$ atau $P\trs R\trs $ merupakan salinan $\BbbR^r$
ke suatu subruang yang memuat $T[\Bbb R^r]$. Semua penyalinan bolak-balik
ini dirancang untuk mengubah $T$ menjadi operator linear
$S:\BbbR^r\to\BbbR^r$ yang dapat kita tangani dengan 263A, dan bagian
(b) dalam bukti memeriksa bahwa kita menyalin ukuran sekaligus struktur
linearnya.

Dalam 265D-265E saya berusaha mengikuti 263C-263D sedekat mungkin.
Sebenarnya hanya diperlukan satu gagasan baru. Ketika $s=r$, kita
mempunyai argumen khusus untuk menunjukkan bahwa
$\mu_r^*\phi[D]\le J\mu_r^*D+\epsilon\mu_r^*D$ (dalam bahasa 263C),
yang berlaku baik $J=0$ maupun tidak. Ketika $s>r$, pendekatan ini gagal,
karena kita tidak lagi dapat menghampiri $\nu_rT[B]$ dengan $\nu_rG$
untuk $G\supseteq T[B]$ terbuka. (Lihat bagian (b-i) bukti 263C.) Karena
itu saya beralih ke argumen lain, yang hanya berlaku ketika $J>0$, dan
akibatnya harus mencari metode terpisah untuk menunjukkan bahwa
$\{\phi(x):x\in D,\,J(x)=0\}$ terabaikan terhadap $\nu_r$. Karena
kita bekerja tanpa batasan pada dimensi $r$, $s$ selain $r\le s$, kita
dapat menggunakan trik menghampiri $\phi:D\to\BbbR^s$ dengan
$\psi_{\eta}:D\to\BbbR^{s+r}$, seperti dalam bagian (d) bukti 265E.

Saya memberikan tiga metode untuk menghitung luas permukaan bola berdimensi
$r$: pendekatan
langsung (265F), metode silinder luar (265Xe), dan teorema integral
berulang yang penting (265G). Dua metode pertama memberikan rumus untuk
luas tudung (265Xb, 265Xe(v)). Metode silinder luar menarik karena
Jacobiannya ternyata $1$, yakni kita mempunyai fungsi \imp\. Saya
mencatat bahwa meskipun telah mengembangkan teknik yang memperbolehkan
domain tidak beraturan, singularitas dalam fungsi $\theta$ pada 265Xe
tetap memaksa saya menangani permukaan bola dalam dua bagian. Teorema 265G adalah
kasus khusus Teorema Koarea ({\smc Evans \& Gariepy 92}, \S3.4;
{\smc Federer 69}, 3.2.12).

Untuk langkah berikutnya dalam teori geometris ukuran pada ruang
Euklides, lihat Bab 47 Jilid 4.

}%end of notes

\frnewpage
